\section{Conclusion}

Completion of this laboratory work has provided practical understanding of binary signal
acquisition and conditioning techniques for embedded sensor systems, and reinforced the
use of FreeRTOS synchronization primitives for safe inter-task communication.

This fifth laboratory work covered the complete sensor pipeline from raw data acquisition
to binary alert generation and structured reporting. The key takeaways are:

\begin{enumerate}
  \item Implementing \textbf{non-blocking sensor acquisition} with the DS18B20 in
    FreeRTOS context -- by using \texttt{setWaitForConversion(false)} and requesting the
    next conversion immediately after reading the previous result, the 94\,ms conversion
    time is hidden inside the 100\,ms task period without blocking the task or the
    scheduler.
  \item Applying \textbf{hysteresis} to convert a continuous analog signal into a
    stable binary alert state -- the two-threshold design (ON $\geq 26\,^\circ$C, OFF
    $\leq 24\,^\circ$C) creates a dead band that eliminates chattering when the
    temperature fluctuates near a single set point.
  \item Combining hysteresis with a \textbf{debounce counter} -- requiring
    \texttt{DEBOUNCE\_COUNT} $= 3$ consecutive confirmations (300\,ms) before accepting
    a state transition prevents brief thermal noise spikes from triggering false alerts,
    which is an essential requirement in industrial sensor conditioning.
  \item Structuring the firmware into three distinct \textbf{FreeRTOS tasks} with
    well-defined roles and priorities: Acquisition (priority 3) runs most frequently and
    ensures fresh data, Conditioning (priority 2) is event-driven via semaphore and
    reacts immediately to new samples, and Reporter (priority 1) provides periodic
    human-readable output without interfering with the control loop.
  \item Using a \textbf{mutex-protected shared data store} (\texttt{SensorData}) as
    the single source of truth for all sensor state, preventing data races when multiple
    tasks read or write temperature, alert state, and debounce counters concurrently.
  \item Organizing the project into \textbf{three hardware-abstraction layers} (MCAL,
    ECAL, SRV) with PlatformIO, keeping all hardware-specific code in the MCAL layer
    and making the service layer fully portable across different MCU platforms.
\end{enumerate}

This laboratory work demonstrates how a small set of classical signal processing
techniques -- hysteresis and debouncing -- combined with FreeRTOS synchronization
primitives, are sufficient to build a robust and noise-resistant binary event detector
from a raw sensor stream.
